% ============================================================================
% UK-RES.tex — category archetype: UK research university
% ----------------------------------------------------------------------------
% LAYER 1 (category archetype) override file for the suite localization model.
% \input AFTER shared/localization-defaults.tex; an institution-level set
% (LAYER 2) is \input AFTER this file to fill the country-specific NAME keys.
%
% ----------------------------------------------------------------------------
% WHICH INSTITUTIONS THIS ARCHETYPE FITS
% ----------------------------------------------------------------------------
% Chartered/statutory research-intensive universities of the United Kingdom
% (England, Scotland, Wales and Northern Ireland) operating the classic
% Council/Senate split: a lay-majority governing body (Council, or Court in
% the older Scottish/post-1992 sense) plus an academic Senate (Academic
% Board), with a Vice-Chancellor as chief executive and academic officer.
% This covers the Russell Group and other research-active pre- and post-1992
% universities connected to Janet via Jisc and bound by the UK
% (post-Brexit) data-and-cyber regime (UK GDPR + DPA 2018 as amended by the
% Data (Use and Access) Act 2025; the NIS Regulations 2018). It does NOT fit
% the federal University of London/collegiate-Oxbridge constituent-college
% governance shape, teaching-led/FE providers, or non-UK institutions; those
% take their own archetype.
%
% ----------------------------------------------------------------------------
% SCOPE OF THIS FILE
% ----------------------------------------------------------------------------
% This is a PARTIAL set. It binds only the STRUCTURAL keys (governance ladder,
% executive leadership, constituent units, NREN tier, identity federation,
% deployment-scale model, data-custody posture, pilot units, operational
% roles) to this archetype's concrete generic value, and points the
% REGULATORY keys at the UK regime worked in shared/regulatory-applicability-
% EN.tex (Section "United Kingdom --- light regime"). It deliberately LEAVES
% the country-specific NAME keys at their generic default --- the specific
% national NREN (Janet/Jisc as the operator name), the data-protection
% authority (the ICO / Information Commission) and the national CSIRT/CERT ---
% so an institution-level set can bind them. Keys not mentioned here keep the
% shared default from localization-defaults.tex.
%
% Requires localization.sty (\setloc / \loc).
% ============================================================================

% --- Subject institution & profile ------------------------------------------
\setloc{institution-type}{a UK research-intensive (typically Russell Group or
  pre-1992 chartered) university}

% --- Institutional governance ladder ----------------------------------------
% Council = lay-majority governing body (corporate/financial/strategic);
% Senate = supreme academic authority. Kept as generic role phrases so the
% archetype fits any UK university regardless of the body's local proper name
% (Council vs Court).
\setloc{governing-board}{the university's lay-majority governing body
  (the Council, or Court) as the supreme corporate, financial and strategic
  authority, working through its CIO and audit/risk committee}
\setloc{academic-governing-bodies}{the academic Senate (Academic Board) and
  its faculty/school boards as the supreme academic authority}
\setloc{executive-leadership}{the Vice-Chancellor (and President/Principal) as
  chief executive and chief academic officer, supported by the senior executive
  team (Provost, Deputy and Pro-Vice-Chancellors, Registrar)}

% --- Constituent units & pilots ---------------------------------------------
\setloc{cs-research-unit}{a Department or School of Computer Science within a
  faculty/college}
\setloc{security-research-centre}{a research-active cybersecurity centre
  (an Academic Centre of Excellence in Cyber Security Research, ACE-CSR)}
\setloc{institution-hpc-facility}{the university's research-computing /
  HPC service (its central advanced-research-computing team)}
\setloc{pilot-entities}{the university's central research-computing / HPC
  service and its strongest-expertise academic units}

% --- Operational roles ------------------------------------------------------
\setloc{internal-emergency-incident-contact}{the university's IT service desk /
  Chief Information Security Officer as the mandatory internal incident-reporting
  route}

% --- Research & education network tier (UK: Janet via Jisc) ------------------
% NREN TIER (structural): the operator NAME (Janet/Jisc) is left to the
% institution-level NAME set; here we bind the structural tier, catalogue,
% backbone scale, federation profile and research-cloud route generically.
\setloc{nren-name}{the UK national research and education network operated by
  the sector's not-for-profit digital body (the Janet Network)}
\setloc{nren-csirt}{the NREN's computer security incident response team
  (the Janet/Jisc CSIRT)}
\setloc{nren-service-catalogue}{the NREN-provided federated-service catalogue
  --- national federated Wi-Fi roaming (eduroam(UK)), a trusted certificate
  service, the UK Access Management Federation single sign-on, security
  operations and managed research-data transfer}
\setloc{research-backbone-scale}{the UK research backbone: a multi-terabit
  national core (Janet, with regional networks aggregating over a thousand
  connected organisations and 18M+ members) connecting at 10--100\,Gbps per
  campus, peering into the European backbone (G\'EANT)}
\setloc{identity-federation-profile}{the UK Access Management Federation
  (UKAMF), interoperating with the academic interfederation (eduGAIN)}
\setloc{nren-research-cloud}{the NREN-brokered research-cloud route ---
  high-bandwidth peering to the major public clouds plus the sector cloud
  framework (OCRE) for EU/UK-compliant providers, rather than an
  NREN-owned cloud}
\setloc{peer-nren-roster}{peer European NRENs (each operating a national
  backbone, a CSIRT and federated identity/storage/IaaS services) coordinated
  through G\'EANT}

% --- High-performance computing (national tier) -----------------------------
\setloc{national-hpc-facility}{a UK national HPC service (the national
  tier-1 supercomputing facility and its regional tier-2 centres)}

% --- Deployment topology & data custody -------------------------------------
% UK post-adequacy posture: keep custody within UK/adequacy-covered space and
% govern transfers under the UK transfer regime (see regulatory keys below).
\setloc{deployment-custody-model}{the university operates, or procures under
  equivalent custody guarantees, a deployment of its core components within the
  UK (or adequacy-covered) jurisdiction}
\setloc{data-custody-sovereignty-profile}{a custody arrangement keeping data
  within the UK / adequacy-covered space and minimising third-country
  disclosure-regime exposure under the UK transfer regime}

% ============================================================================
% REGULATORY KEYS --- UK regime (post-Brexit, light regime)
% ----------------------------------------------------------------------------
% These reference the worked UK instantiation in
% shared/regulatory-applicability-EN.tex (Section "United Kingdom --- light
% regime", Table tab:reg-uk). They bind the REGIME (instruments / scope), not
% the named authority: the ICO/Information Commission, the specific UK
% data-protection authority and national CSIRT NAMES are left to the
% institution-level NAME set.
% ============================================================================
\setloc{regulatory-regime}{binding UK law (retained-and-amended EU-derived
  instruments now diverging via domestic reform)}
\setloc{data-protection-regime}{the UK GDPR with the Data Protection Act 2018,
  as amended by the Data (Use and Access) Act 2025}
\setloc{cybersecurity-baseline-regime}{the UK cyber-resilience framework --- the
  NIS Regulations 2018 (SI 2018/506) with the NCSC Cyber Assessment Framework,
  ahead of the forthcoming Cyber Security and Resilience Bill}
\setloc{eid-trust-services-regime}{the UK eIDAS regime (retained
  Reg.~(EU)~910/2014 as amended; SI~2016/696)}
\setloc{cross-border-transfer-regime}{the UK GDPR Chapter~V transfer regime ---
  adequacy where it exists, else the ICO International Data Transfer Agreement
  (IDTA) or the UK Addendum to the EU SCCs, with a transfer-risk assessment}
\setloc{ai-regulation}{the UK's principles-based, regulator-led AI-governance
  approach (sectoral regulators applying the cross-sector AI principles, no
  single omnibus AI statute)}
\setloc{essential-entity-classification-basis}{operators within scope of the
  NIS Regulations 2018, with sector competent authorities (CSR Bill expanding
  scope to managed service providers, not yet in force)}

\endinput
